\documentclass[12pt]{article}
\usepackage[utf8]{inputenc}
\usepackage[T1]{fontenc}
\usepackage{hyperref}
\title{Karuwali Language Resources for K-12 Curriculum Development}
\author{Compiled by Claude Sonnet 5.5 (claude-sonnet-5-5) for Andr\'as Kornai}
\date{}
\begin{document}
\maketitle

\section{Introduction}
Karuwali (also Garuwali) is a dormant Pama-Nyungan language of far-western Queensland. It is associated with the Diamantina Shire area, including Betoota and Haddon Corner. Norman Tindale estimated its territory at about 31,000 square kilometres, from near Farrars Creek south to Betoota and Haddon Corner. The search results gave little beyond encyclopedic and archival listings.

\section{Primary Educational Resources}
No primers, textbooks or school programmes for Karuwali were found.

\section{Linguistic and Reference Materials}
\begin{itemize}
\item Wikipedia, Karuwali: \url{https://en.wikipedia.org/wiki/Karuwali}
\item South Australian Museum archive, Karuwali (QLD): \url{https://www.samuseum.sa.gov.au/collection/archives/language_groups/karuwali}
\item Wikipedia, Mithaka language, a neighbouring language that places Karuwali in the Pama-Nyungan family: \url{https://en.wikipedia.org/wiki/Mithaka_language}
\end{itemize}
A search summary said AIATSIS holds a bibliography of Garuwali materials. The exact URL was not seen, so none is given here.

\section{Religious and Community Resources}
No Karuwali community organisation or language programme was found.

\section{Digital and Audio Resources}
No audio, video or software was found.

\section{Curriculum Development Context}
Material for a curriculum is very thin in the results. The first step would be to approach AIATSIS and the Diamantina Shire community and Traditional Owners to locate archival word lists and recordings. Regional context is in \url{https://en.wikipedia.org/wiki/Shire_of_Diamantina}.

\end{document}